\documentclass{article}
\usepackage{pathfinder-readable}

\title{A Carbon Forecast in a Resource Allocation Game: A Restricted Result}

\begin{document}
\maketitle

\begin{abstract}
Paper Q designs payments for strategic users of a shared resource, while paper P forecasts the carbon
intensity of electricity and schedules movable demand. The thread asked whether Q's game could put P's
forecast to work in decentralised scheduling. The peers found faults in Q's claims about unique prices
and voluntary participation, then checked a corrected payment that gives a unique forecast-optimal
allocation, a balanced budget and voluntary participation for optional convex loads. That result does not
cover P's fixed workload or mobile-storage routes, and a forecast-optimal allocation need not reduce
actual generation emissions. The thread ended as a draft because the link from forecast carbon accounting
to physical emissions remains unproved.
\end{abstract}

\section{Introduction}

Q asks how a manager can share limited resources among users who know their own benefits. Each user
requests an amount and proposes a price. The manager uses a quadratic payment to make the users' choices
serve a common objective. Q states that its rule has a unique Nash equilibrium, balances payments, gives
users a reason to participate and supports a learning algorithm \cite{Q}.

P forecasts the carbon intensity of electricity at each network location and time. It then uses those
day-ahead figures to schedule flexible demand, including data-centre work and mobile energy storage. Its
scheduling problem includes workload that must be served and storage routes that require yes-or-no travel
decisions. Its reported test uses a modified IEEE 33-bus system \cite{P}.

The proposed connection was to put P's forecast into a Q-style payment. Users could then choose loads in
response to a carbon charge while the resulting allocation followed the forecast-based plan. The peers
asked two prior questions: whether Q's guarantees held for its displayed payment, and whether P's
scheduling and carbon measure fit the assumptions needed for those guarantees. Their answers narrowed the
connection to an optional, continuous load market. Here ``optional'' means that each user can choose zero
load.

\section{What was done}

The peers first compared the two models. Q's local choices are continuous and convex, and users share
componentwise upper capacity limits. P has a continuous data-centre workload component, but its total
workload is fixed and its mobile-storage model has binary route choices. The peers therefore did not
treat P's full programme as an instance of Q's game. They also noticed that P calculates nodal intensity
using generation, power flows and storage discharge, then schedules against a forecast held fixed during
the decision \cite{Q,P}.

Next, Emmy checked Q's uniqueness conditions. One condition requires a strictly negative quadratic form
for all nonzero directions; another requires certain price-block coefficients to sum to zero. In a
direction where all users change their price by the same amount and no allocation changes, the latter
condition makes that quadratic form zero. The two conditions cannot both hold. Emmy then gave a two-user
example in which Q's displayed payment has many equilibrium prices but only one efficient allocation. Ada
independently checked participation and found an equilibrium at which one user's utility is negative,
despite efficient allocation and a balanced budget. Both checks are recorded in their derivations and the
ledger.

The peers then worked back from Q's displayed payment, rather than its claimed uniqueness certificate.
Emmy derived its price best responses and the allocation result they imply. Ada found a payment
adjustment based only on other users' actions, so it changes no user's best response. She combined it
with a carbon charge based on P's frozen forecast and a rebate based on other users' loads. Emmy checked
the adjustment and the two forecast-error bounds; Ada added the condition needed to speak about physical
generation emissions. The remaining objections concerned P's required workload, its route decisions and
the difference between average intensity and emissions caused by an added load.

\section{Results for optional loads}

Here is the precise scope of the repaired result. There are \(N\geq2\) users and \(K\) load coordinates,
such as bus and hour pairs. User \(n\) chooses a nonnegative vector \(x_n\) from a compact convex set
\(\mathcal X_n\) that contains zero. Its benefit \(V_n(x_n)\) is strongly concave and \(V_n(0)=0\). The
users' total load \(X=\sum_n x_n\) cannot exceed a nonnegative capacity vector \(c\), coordinate by
coordinate. The positive constant \(\alpha\) scales Q's payment. A nonnegative vector \(\hat e\) gives
P-style forecast intensities, held fixed while users choose. A nonnegative constant \(\beta\) weights
carbon in the objective. Assume that the resulting planner problem has an optimum and a capacity
multiplier. The planner maximises
\[
F_{\hat e}(x)=\sum_{n=1}^{N}V_n(x_n)-\beta\hat e^{\top}X
\quad\text{subject to}\quad x_n\in\mathcal X_n,\qquad X\leq c.
\]

The central check is short. Let \(p_n^k\geq0\) be user \(n\)'s proposed price for coordinate \(k\), let
\(S^k=X^k-c^k\) be unused capacity with its sign reversed, and let \(P^k=\sum_n p_n^k\). Q's displayed
rule gives the price best response
\[
p_n^k=\left[\frac{P^k-p_n^k+S^k}{N-1}\right]_+,
\]
where brackets mean the maximum of the enclosed value and zero. If any price for a coordinate is
positive, all users' prices for it must be equal and its capacity must be exactly filled. If all prices
are zero, the capacity cannot be exceeded. Thus every equilibrium has a common nonnegative price vector
\(p\), feasible loads and a zero product between each price and its capacity slack. At that common price,
a user's marginal payment for load is \(\alpha p\). The users' allocation conditions are therefore the
planner's conditions, with capacity multiplier \(\alpha p\). Conversely, a planner optimum and multiplier
give an equilibrium. Strong concavity makes the allocation unique, although the price need not be unique.
Emmy's derivation gives the full best-response and converse checks \cite{Q}.

Q's raw payment, denoted \(t_n^Q\) for user \(n\), charges too much in the participation example below.
Ada's correction adds a transfer \(h_n\) that depends only on the other users' prices and loads. Writing
\(\bar p_{-n}=\sum_{m\ne n}p_m\) and \(\bar x_{-n}=\sum_{m\ne n}x_m\), it is
\[
h_n=\frac{\alpha N}{(N-1)^2}
\left(\frac{\bar p_{-n}}{N-1}\right)^{\top}
\left(\bar x_{-n}-\frac{N-1}{N}c\right).
\]
The forecast charge with its rebate is
\[
b_n=\beta\hat e^{\top}x_n-
\frac{\beta}{N-1}\hat e^{\top}\bar x_{-n}.
\]
The repaired payment is \(t_n=t_n^Q+h_n+b_n\). Neither \(h_n\) nor the rebate part of \(b_n\) depends on
user \(n\)'s action. Its choices are therefore those under Q's game with benefit \(V_n(x_n)-\beta\hat
e^{\top}x_n\). At equilibrium, direct substitution reduces \(t_n^Q+h_n\) to \(\alpha p^{\top}(x_n-c/N)\).
These terms sum to zero because a positive price implies filled capacity; the \(b_n\) terms sum to zero
at every profile. The repaired game is budget balanced at equilibrium.

For participation, zero load is a feasible choice. The user's allocation best response implies that its
benefit, less the forecast carbon charge and \(\alpha p^{\top}x_n\), is at least zero. Its equilibrium
utility also includes the nonnegative amounts \(\alpha p^{\top}c/N\) and \(\beta\hat e^{\top}\bar
x_{-n}/(N-1)\). It is therefore at least zero relative to the zero-service option. Ada's derivation and
Emmy's independent check support the payment and participation calculations.

There is a conditional forecast guarantee. Let \(\mathcal F\) be the feasible allocations and let
\(e(x)\) be the intensity that would result if allocation \(x\) were chosen. If every feasible
counterfactual obeys \(\lVert e(x)-\hat e\rVert_\infty\leq\varepsilon\), the loss of the forecast-optimal
allocation in the objective that charges \(\beta e(x)^{\top}X\) is at most \(2\beta\varepsilon\lVert
c\rVert_1\). Here \(\varepsilon\) is a uniform forecast-error limit, \(\lVert\cdot\rVert_\infty\) is the
largest coordinate magnitude, and \(\lVert c\rVert_1\) is the sum of capacities. The proof compares the
true and forecast objectives at both the forecast choice and the true optimum. If true intensity is a
fixed vector, the tighter limit is \(\beta\varepsilon\lVert c\rVert_1\). Both limits concern the carbon
attributed to the chosen load.

To bound loss in actual generation emissions, Ada required one more premise. Let \(C(x)\) be generation
emissions under \(x\), let \(C_0\) be a fixed baseline, and suppose \(\lvert
C(x)-C_0-e(x)^{\top}X\rvert\leq\delta\) for every feasible \(x\). Then loss in the objective using
\(C(x)-C_0\) is at most \(2\beta(\varepsilon\lVert c\rVert_1+\delta)\). The nonnegative number \(\delta\)
limits the gap between attributed load emissions and changed generation emissions. P's reported forecast
tests establish neither of these uniform premises \cite{P}.

\section{What did not hold}

Q's strict uniqueness certificate fails the common-price direction test described above. Full equilibrium
uniqueness also fails directly. Give each of two users a load in \([0,1]\) and benefit \(3x-x^2/2\). Set
capacity to \(2\) and \(\alpha=1\). Every profile with both loads equal to \(1\) and both prices equal to
any value from \(0\) to \(2\) is an equilibrium. The allocation is unique. This is a failure of unique
prices, not of the allocation argument \cite{Q}.

Q's displayed payment also fails its stated individual-rationality claim. In Ada's two-user example,
capacity and \(\alpha\) are both \(1\); the users' benefits are \(3x-x^2/2\) and \(2x-x^2/2\), with each
load in \([0,1]\). Loads \((1,0)\) and prices \((2,2)\) form an equilibrium. Q's payments are \((3,-3)\),
leaving the first user with utility \(-1/2\). The repaired transfer removes this failure for the
optional-load setting; it does not establish participation for P's workload that must be served.

The carbon objection remains. In Emmy's two-location thought experiment, location A's zero-carbon supply
already meets its baseline load and is at capacity. An extra unit at A uses backup generation that emits
one unit. Location B can supply that unit with one-half unit of emissions. Baseline average intensity
ranks A as cleaner, so a scheduler using it sends the movable unit to A, although B causes fewer added
emissions. This is a constructed check, not a result from P's test network. It shows why even an exact
baseline intensity forecast need not rank physical responses correctly.

The peers also checked prior work. Chen and Zhao already claim social optimality, budget balance and
individual rationality in a joint electricity-carbon pricing mechanism \cite{ChenZhao}. The ledger
contains no exhaustive review. A balanced carbon charge alone therefore cannot carry a novelty claim
here. Nor did the peers establish stochastic learning convergence for the repaired game or a mechanism
for P's binary mobile-storage routes.

\section{Why the thread stopped}

The verifier left the thread at DRAFT after one round. It accepted the narrower, checked result: a
corrected payment implements the unique forecast-conditioned allocation with balanced payments and
voluntary participation for optional convex loads. It also accepted the note's limits: P's fixed-workload
and binary-routing programme is outside that result, and no physical emissions reduction follows from the
forecast alone. The smallest useful restart is one explicit two-location movable-load dispatch: calculate
generation emissions \(C(x)\) and counterfactual intensity \(e(x)\) for its feasible choices, then test
the two uniform bounds above. That would show whether the accounting guarantee says anything about
physical emissions in that instance.

\bibliographystyle{plainurl}
\bibliography{references}

\end{document}
